\section{構造的解釈}
\label{sec:structural}

\subsection{構造クラスの候補}
VED$\times$IFGT枠組みは、以下の構造クラスの候補と認知的解釈を動機づける。
表はモデルに基づく分類を整理するものであり、形式的な存在定理や分類定理ではない。
詳細な解の解析と独立した認知的対応は今後の課題である。

\begin{center}
\renewcommand{\arraystretch}{1.5}
\begin{tabular}{p{0.29\textwidth}p{0.31\textwidth}p{0.30\textwidth}}
\toprule
\textbf{候補クラス} & \textbf{場の構造} & \textbf{認知的解釈} \\
\midrule
滑らかな勾配領域 & $\nabla\Phi$が緩やかに変化 & 背景的な認知的傾向 \\
浅い準安定井戸 & $\nabla\Phi\approx 0$、深さが小さい & 暫定的または不安定な概念 \\
深い安定井戸 & $\nabla\Phi\approx 0$、深さが大きい & 定着した概念や信念 \\
リミットサイクル軌道 & $\JI$が沈殿せず循環 & 反復する思考パターン \\
障壁領域 & $|\nabla\Phi|$が大きい & 認知的な行き詰まりや抵抗 \\
尾根構造 & $\Phi$の局所最大 & 概念境界やカテゴリの縁 \\
\bottomrule
\end{tabular}
\end{center}

\subsection{深さ、幅、頑健性}
三つの幾何的パラメータが概念アトラクタを特徴づける。

\textbf{深さ}（井戸内の$-\min\Phi$）には、新たな入力による変形への抵抗との対応を提案する。
深い井戸を抵抗の強い概念と、浅い井戸を変位しやすい概念と比較する。

\textbf{幅}（閾値未満の領域の空間的広がり）は、文脈の範囲と比較する。
広い井戸を適用範囲の広さと、狭い井戸を文脈依存性の高い概念と比較する。

\textbf{頑健性}（極小点における$\Phi$の曲率）は、精密さの比較指標として提案する。
高い曲率を鋭い概念と、低い曲率を曖昧または多義的な概念と定性的に関連づける。
これらの意味的属性は、いずれも幾何的記述だけでは較正されない。

これらの幾何的属性は、固定された内在的性質として扱わない。
式\eqref{eq:I_evolution}--\eqref{eq:sigma_barrier}の関係のもとで、$\Phi$とともに発展するものとしてモデル化する。

\subsection{概念境界}
$\Phi$の尾根構造、すなわち隣接する井戸を隔てる局所的に高いポテンシャルの領域は、この記述では概念間の区別である概念境界を表す。
境界は鋭い線ではなく、通過コストの高い広がりを持った領域である。

定性的な仮説は、反復的な言語的・経験的摂動のもとで概念境界が移動しうる、というものである。
尾根近傍で結合するトークンは、尾根を強化して区別を深めたり、侵食して隣接概念を融合させたりすると提案する。
これは意味ドリフトと概念変化の動的記述の候補であり、確立された認知機構や意味の対応写像ではない。

\subsection{語彙の選択と分布的記述}
安定化したトークン$\{\ell_i\}$の全体集合を、この枠組みでは語彙と呼ぶ。
以下の表式は安定性によって集合を選ぶものである。
語彙が場に分布するという広い記述によって、その集合が確率分布になるわけではない。
\begin{equation}
  \mathcal{V} = \bigl\{\ell_i \;\big|\; S(\ell_i) > \epsilon \bigr\},
  \label{eq:vocabulary}
\end{equation}
ここで$S(\ell_i)$は、再利用頻度、持続時間、盆地の深さに依存することを意図した安定性汎関数であり、$\epsilon$は系に依存する閾値である。
ここでは、この汎関数と閾値に抽出規則や較正規則を与えない。

語彙は、使用によって再形成される、場に関連した組織化として解釈する。
トークン頻度には盆地のアクセス可能性との関係を提案するが、意味的重要性の尺度とはしない。
曖昧で文脈依存的な区別は、この解釈的描像の一部であり、トークンの操作的基準ではない。
